\section{逐区需求与结果表口径}\label{app:outputs}
\subsection{逐服务区需求统计}
表\ref{tab:demand-details}依据逐箱清单汇总，硬时限箱按医疗与首批要求的并集计算。正文中的总质量、体积和31箱硬时限均可由该表回算。
\begin{longtable}{lrrrr}
\caption{逐服务区需求统计}\label{tab:demand-details}\\
\toprule 服务区 & 箱数 & 质量/kg & 体积/m$^3$ & 硬时限箱数\\\midrule\endfirsthead
\toprule 服务区 & 箱数 & 质量/kg & 体积/m$^3$ & 硬时限箱数\\\midrule\endhead
S001 & 15 & 154.00 & 0.394 & 3 \\
S002 & 8 & 81.00 & 0.211 & 2 \\
S003 & 8 & 81.00 & 0.211 & 2 \\
S004 & 6 & 59.00 & 0.156 & 2 \\
S005 & 6 & 59.00 & 0.156 & 2 \\
S006 & 6 & 59.00 & 0.156 & 2 \\
S007 & 5 & 45.00 & 0.129 & 2 \\
S008 & 5 & 45.00 & 0.129 & 2 \\
S009 & 3 & 25.00 & 0.067 & 2 \\
S010 & 3 & 25.00 & 0.067 & 2 \\
S011 & 3 & 25.00 & 0.067 & 2 \\
S012 & 3 & 25.00 & 0.067 & 2 \\
S013 & 3 & 25.00 & 0.067 & 2 \\
S014 & 3 & 25.00 & 0.067 & 2 \\
S015 & 3 & 25.00 & 0.067 & 2 \\
合计 & 80 & 758.00 & 2.011 & 31 \\
\bottomrule\end{longtable}


\subsection{逐箱、逐架次和通信分段的输出}
\begin{table}[H]\centering
\caption{结果记录的对象与关键字段}\label{tab:output-fields}
\begin{tabularx}{\linewidth}{p{2.6cm}Y}\toprule
记录对象 & 关键字段与解释\\\midrule
单点批次 & 服务区、货箱集合、机型、质量、体积、能耗、完整作业时长\\
运输架次 & 架次编号、访问顺序、机型、实体机、电池、准备开始、起飞、返航、能耗与SOC\\
逐箱交付 & 货箱编号、所属架次、服务区、站内交接序号、交付完成、硬截止或期望时刻\\
能源周转 & 能源编号、任务区间、返航SOC、充电开始和结束、下次分配任务\\
中继架次 & 悬停经纬度与海拔、离地高度、实体机、能源组件、抵达、建链、服务结束和返航\\
通信分段 & 运输架次、飞行或交接阶段、区间边界、直连或中继状态、实际中继编号、可用性依据\\
分区配置 & 服务区集合、关联运输及中继任务、分类型资源数、全局缺口、相对集中执行冗余与工作量\\
\bottomrule\end{tabularx}\end{table}

表\ref{tab:output-fields}中的“开始”统一指准备开始，“交付”指单箱交接完成，“任务完成”指最后返航。电池及能源组件的充电结束只用于下一任务可用性，不加入已完成任务的完工时间。各结果表由同一计算程序导出，字段口径见表\ref{tab:output-fields}。
